\chapter{Tình cảm và lòng người}
\markboth{Tình cảm và lòng người}{Love and Human Hearts}

\section{Một người chỉ quý sự dịu dàng sau khi đã quen với lạnh nhạt.}
\begin{truyen}{Một người chỉ quý sự dịu dàng sau khi đã quen với lạnh nhạt.}{Gentleness is a rare wealth.}
\chuhoa{S}{au nhiều mối quan hệ khiến lòng mệt, cô mới gặp một người biết lắng nghe, không cợt nhả nỗi đau, không lấy cơn nóng giận làm quyền lực. Khi ấy cô mới hiểu dịu dàng không hề yếu; nó là một phẩm chất rất hiếm.}
\textit{(After relationships that exhausted her, she met someone who listened, did not joke about pain, and did not use anger as power. Then she understood that gentleness is not weak; it is rare.)}
\end{truyen}
\begin{baihoc}
Con người thường hiểu muộn rằng được đối xử tử tế là một điều quý, không phải chuyện hiển nhiên.
\textit{(People often learn too late that being treated kindly is precious, not automatic.)}
\end{baihoc}
\begin{ghinhoanh}
\textbf{Short English to remember:}

Gentleness is a rare wealth.

Kind love is not a small thing.
\end{ghinhoanh}
\begin{tuhocnhanh}
1. Em có đang xem nhẹ sự dịu dàng của ai đó dành cho mình không? \textit{(Are you undervaluing someone's gentleness toward you?)}

2. Em có đang dịu dàng đủ với người mình thương không? \textit{(Are you gentle enough with the people you love?)}
\end{tuhocnhanh}
\ngancach

\section{Một người ở trong mối quan hệ chỉ vì sợ cô đơn.}
\begin{truyen}{Một người ở trong mối quan hệ chỉ vì sợ cô đơn.}{Fear is a poor reason to stay.}
\chuhoa{A}{nh biết mình không còn được tôn trọng, nhưng nghĩ chia tay sẽ trống trải hơn. Về sau, anh mới nhận ra ở lại trong một nơi làm mình nhỏ đi cũng là một kiểu cô đơn, chỉ kéo dài hơn và đau hơn.}
\textit{(He knew he was no longer respected, but thought leaving would be lonelier. Later he realized that staying in a place that makes you smaller is also a kind of loneliness, only longer and more painful.)}
\end{truyen}
\begin{baihoc}
Không phải cứ có người bên cạnh là lòng sẽ bớt trống; có những mối quan hệ chỉ làm trống thêm.
\textit{(Having someone beside you does not always ease emptiness; some relationships deepen it.)}
\end{baihoc}
\begin{ghinhoanh}
\textbf{Short English to remember:}

Fear is a poor reason to stay.

Loneliness inside love still hurts.
\end{ghinhoanh}
\begin{tuhocnhanh}
1. Em có đang ở lại đâu đó chỉ vì sợ một mình không? \textit{(Are you staying somewhere only because you fear being alone?)}

2. Em cần điều gì hơn: người ở cạnh hay sự tôn trọng? \textit{(What do you need more: company or respect?)}
\end{tuhocnhanh}
\ngancach

\section{Một người chỉ hiểu giá trị của tin nhắn ngắn khi nó ngừng đến.}
\begin{truyen}{Một người chỉ hiểu giá trị của tin nhắn ngắn khi nó ngừng đến.}{Small care becomes visible in absence.}
\chuhoa{M}{ỗi sáng có người nhắc ăn sáng, hỏi ngủ ngon không, dặn đi đường cẩn thận. Khi những dòng ngắn ấy biến mất, cô mới hiểu sự quan tâm không luôn mang hình thức lớn; đôi khi nó nằm trong những điều lặp lại rất nhỏ.}
\textit{(Each morning someone reminded her to eat, asked if she slept well, and told her to travel safely. When those short messages disappeared, she understood that care is not always dramatic; sometimes it lives in small repetitions.)}
\end{truyen}
\begin{baihoc}
Điều bền bỉ thường đáng quý hơn điều rực rỡ; ta chỉ hay hiểu ra khi nó không còn nữa.
\textit{(What is steady is often more valuable than what is dramatic; we just tend to realize it when it is gone.)}
\end{baihoc}
\begin{ghinhoanh}
\textbf{Short English to remember:}

Small care becomes visible in absence.

Do not wait for loss to notice love.
\end{ghinhoanh}
\begin{tuhocnhanh}
1. Có sự quan tâm nhỏ nào em đã quen đến mức không còn thấy quý? \textit{(What small act of care have you become too used to appreciate?)}

2. Em có thể đáp lại sự quan tâm đó ra sao? \textit{(How can you respond to that care?)}
\end{tuhocnhanh}
\ngancach

\section{Người ta nói lời nặng nhất với người thân nhất.}
\begin{truyen}{Người ta nói lời nặng nhất với người thân nhất.}{Closeness is no license for cruelty.}
\chuhoa{V}{ì tin rằng người kia sẽ hiểu, người ta dễ buông lời nặng hơn, nóng hơn, bất cẩn hơn. Đến khi khoảng cách xuất hiện, ta mới thấy sự gần gũi không cho phép mình thiếu tử tế.}
\textit{(Because they assume the other person will understand, people often speak more harshly, more angrily, and more carelessly to those closest to them. When distance appears, they realize closeness does not permit cruelty.)}
\end{truyen}
\begin{baihoc}
Người thân nhất là người càng đáng được nghe lời mềm nhất.
\textit{(The people closest to us deserve the gentlest words most.)}
\end{baihoc}
\begin{ghinhoanh}
\textbf{Short English to remember:}

Closeness is no license for cruelty.

Be soft where love is near.
\end{ghinhoanh}
\begin{tuhocnhanh}
1. Em có đang dễ lời nặng với người gần mình nhất không? \textit{(Are you speaking most harshly to those closest to you?)}

2. Câu nào em nên nói mềm lại từ hôm nay? \textit{(Which kind of sentence should you soften from today onward?)}
\end{tuhocnhanh}
\ngancach

\section{Một người hiểu quá muộn rằng được yêu không phải là được chiều hết.}
\begin{truyen}{Một người hiểu quá muộn rằng được yêu không phải là được chiều hết.}{Love is not endless indulgence.}
\chuhoa{C}{ô từng nghĩ ai yêu mình sẽ luôn nhường, luôn chịu, luôn quay lại. Khi người kia mệt mỏi và dừng lại thật, cô mới hiểu tình yêu cũng có giới hạn nếu một bên chỉ nhận mà không chịu trưởng thành.}
\textit{(She thought that anyone who loved her would always give in, endure, and return. When the other person truly grew tired and stopped, she realized that love also has limits when one side only takes and never grows.)}
\end{truyen}
\begin{baihoc}
Tình yêu không phải chiếc chăn che hết mọi thiếu sót; nó cần được giữ bằng sự tử tế và tự sửa mình.
\textit{(Love is not a blanket that covers every flaw; it must be sustained by kindness and self-correction.)}
\end{baihoc}
\begin{ghinhoanh}
\textbf{Short English to remember:}

Love is not endless indulgence.

Being loved is not a right to be careless.
\end{ghinhoanh}
\begin{tuhocnhanh}
1. Em có đang xem tình yêu như điều hiển nhiên không? \textit{(Are you treating love as something guaranteed?)}

2. Em cần trưởng thành ở điểm nào để giữ một mối quan hệ tốt hơn? \textit{(In what area do you need to mature to sustain a relationship better?)}
\end{tuhocnhanh}
\ngancach